\chapter{玩梗接入产业：授权体系、二创周边与文化收编}\label{ch:09}

\begin{nlkeypoints}
  \begin{enumerate}[label=(\arabic*),leftmargin=2.4em]
    \item 正典与变种走上两条经济轨道：
          官方奶龙走\keyw{“内容 + 授权”双轮驱动}
          （200+ 品牌、2000+ SKU、57 国，\conf{A/B} 级），
          奶蛙走\keyw{零授权的可传播性经济}——
          内容层无人收费、形式层无法收费，
          被商业捕捉的是立场的周边化（立牌、玩偶）
          （本书分析，\conf{M} 级）。
    \item 2026-03 关评拉黑事件应读作“不彻底的抑制”：
          官方控制的是\keyw{可执行对象}（平台评论区、账号），
          而“捧腹大笑”作为\keyw{可传播性单位}
          在平台空间之外继续流通——
          \keyw{法律层面}与\keyw{意义层面}的
          正典—变种关系由此分离（\conf{C} 级事件 + \conf{M} 级分析）。
    \item Hebdige 的“收编”在奶蛙案例中呈现为双向运动：
          商业收编迷因（周边商品化），
          迷因也反向收编商业（品牌借梗营销）；
          第一版教材记录的淘宝周边与
          \en{“奶龙正在成为艺术品”}评述
          是这场双向运动的两端（均 \conf{C} 级）。
  \end{enumerate}
\end{nlkeypoints}

第一章留下过一个张力表述：一只“不属于任何官方”的青蛙，
寄生在一个有专业法务团队的正版 IP 旁边生长。
本章把这个张力展开为完整的经济与法律分析。
先立底图（融合文化与可传播性），再看两条轨道各自的
收益结构，然后处理 2026-03 事件的法律—意义双重读法，
最后讨论收编与艺术化的双向运动。
本章继续遵守全书纪律：产业数字全部来自 \conf{A/B} 级
公开报道，二创周边现象来自 \conf{C} 级第一版教材记录，
凡属本书的经济学解释一律标注 \conf{M} 级。

\section{底图：融合文化与可传播性}\label{sec:09:base}

第二章已给出媒介结构的位置：正典奶龙由制作方控制分发，
奶蛙的流通则完全由二创者的转发、二创与线下搬运驱动。
本节把这张底图补上经济维度。Jenkins 等人提出，
在融合文化环境下，内容的价值不取决于制作方的
“黏性”投放，而取决于受众的\keyw{可传播性}
（\en{spreadability}）行为——受众愿意并且能够
携带内容去哪里，内容的价值边界就在哪里
\cite{jenkins2013spreadable}。

用这个视角回看第四章的一个悬置问题：
“结尾集体大笑”为什么值得当作\keyw{可传播性单位}分析
（\ref{sec:04:audio}）。答案现在可以给出：

\begin{proposition}[可传播性单位（教学命题）]\label{prop:09:unit}
  一个文化构件的可传播性，取决于它在搬运中
  \keyw{不依赖原始语境}的能力。“结尾集体大笑”
  是奶蛙语料中可传播性最高的构件：
  它不依赖任何具体压力桥段（内容可换）、
  不依赖原始上传者（署名可弃）、
  不依赖特定平台（载体可换），
  只依赖观看者的“接梗”能力（社群在场）。
\end{proposition}

这条命题（\conf{M} 级）解释了第九章的核心张力：
\keyw{可传播性最强的构件，恰恰是法律上最无法据为己有的构件}。
笑声是录音制品，套路是思想而非表达，梗图与视频
分散在无数匿名账号——官方奶龙可以把形象授权给 200 家品牌，
却无法对“结尾一起笑”这个形式收取任何费用。
可传播性与可产权化在奶蛙案例中呈反向关系，
这是迷因经济的第一个结构性事实。

\section{正典轨道：“内容 + 授权”双轮驱动}\label{sec:09:canon}

第三章已经给出正典奶龙的产业量化口径
（表 \ref{tab:03:numbers}），本节不再复述数字，
只分析其经济结构的三个特征。

\begin{enumerate}[label=\arabic*.]
  \item \textbf{收入的双轮结构}。谢承恩自述：
        2020 年初“商业化也还没起步”，其后
        “从广告合作起步，逐步拓展到商业授权，
        积极探索‘内容 + 授权’双轮驱动的模式”
        \conf{A}\cite{jiang2026nailong}。
        内容轮（平台分账、播放流量）与授权轮
        （品牌合作、衍生品）的比重变化
        没有公开数据，但 200+ 授权品牌、2000+ SKU、
        57 国销售、25 万元政府扶持奖励
        \conf{A/B}\cite{jiang2026nailong,chinadaily2026six}
        说明授权轮已是产业主体。
  \item \textbf{授权的本质是“形象租赁”}。品牌方购买的
        不是剧集内容而是\keyw{形象的情感附加值}——
        “治愈、欢乐”的定位\conf{A}\cite{jiang2026nailong}
        使奶龙成为低风险的情绪符号，
        适配茶饮、乐园、大屏等家庭场景
        （第三章案例 3-1）。
  \item \textbf{维权是授权的前置条件}。法务团队、
        快速维权机制、海外维权协助
        \conf{A}\cite{jiang2026nailong}
        不是授权体系的外挂，而是其成立条件：
        \keyw{只有能把形象圈起来的权利人，才能把形象租出去}。
        维权强度因此可以读作授权体系健康状况的代理指标。
\end{enumerate}

\begin{nlcase}[案例 9-1：表情包的 80 亿次]
  中国日报网报道给出一个耐人寻味的口径：
  奶龙微信表情包累计发送量突破 80 亿次
  （截至 2025 年底，\conf{B}\cite{chinadaily2026six}）。
  表情包位于授权体系与迷因流通的\keyw{交界地带}：
  它是官方发布、免费分发、用户自发使用的——
  形象受控，用途不受控。80 亿次发送既是
  授权体系对迷因式流通的\keyw{主动模仿}
  （官方把形象做成“可传播性单位”投放），
  也说明了授权体系的天花板：
  用户对表情包的每一次再加工、再配文，
  都在官方控制之外。对照 2026-03 对奶蛙的拉黑可见，
  官方并不反对迷因式流通\keyw{本身}，
  反对的是\keyw{失控的流通}——
  这一区分是理解两条轨道摩擦的关键。
\end{nlcase}

\section{变种轨道：零授权的可传播性经济}\label{sec:09:variant}

第一版教材第八部分记录了奶蛙周边的民间售卖：
“淘宝售卖的奶脚抱枕、奶蛙玩偶、亚克力”，
“2m 奶蛙人形立排和奶蛙六合一战神金刚”，
并评论“紧随抽象潮流的商家们也纷纷上架奶蛙周边”
\conf{C}（第一版教材口径；具体商品图本书不复制，
权利状态不明）。

这条轨道的经济结构与正典轨道逐项相反：

\begin{table}[htbp]\centering
  \caption{两条轨道的经济结构对照（本书整理；
          官方侧为 \conf{A/B} 级口径，民间侧为 \conf{C} 级记录，
          “对照”列为 \conf{M} 级分析）}
  \label{tab:09:two}
  \small
  \begin{tabularx}{\textwidth}{@{}lXX@{}}
    \toprule
    维度 & 正典轨道（官方奶龙） & 变种轨道（奶蛙周边） \\ \midrule
    权利状态 & 权利清晰、集中持有 &
      权利模糊：形象未经授权，署名权在匿名创作者 \\
    收益形式 & 授权费、衍生品分成、平台分账 &
      商家现货差价；创作者基本不获酬 \\
    质量控制 & 官方审校、SKU 管理 &
      无控制：从亚克力到 2 米立排质量参差 \\
    法律风险 & 风险在侵权方 &
      风险在售卖方：随时可能被维权 \\
    传播逻辑 & 黏性投放 + 节点营销 &
      可传播性搬运 + 刷屏 \\ \bottomrule
  \end{tabularx}
\end{table}

表中最要紧的一行是\keyw{创作者不获酬}。
奶蛙的爆火使形象产生了真实的商业价值
（周边有市场、立牌有仪式功能），
但按第四章的处理，考据指向的创建者是
\keyw{权利主体}而不只是证据节点（规则 E3）；
社群传播中的每一次免费使用，
都在没有报酬通道的情况下增殖着这一形象的价值。
这不是本书对任何一方行为的法律指控
——恰恰相反，报酬通道的缺失是
\keyw{结构性的}：匿名使授权无从谈起，
授权的缺席使商业化只能在灰色地带进行。
第十二章把这一问题归入“研究伦理三条线”的第三线。

\begin{nlcritique}[批判性反思：谁在赚迷因的钱]
  民间周边经济中存在一个容易忽视的不对称：
  \keyw{承担形象诞生风险的人（创作者）与
  收获售卖利润的人（商家）不是同一批人}。
  第一版教材以轻松笔调记录周边“壮大”，
  本书保留其史料价值，但必须补上另一面：
  无授权售卖使商家处于持续的法律风险中，
  也使消费者买到的是权利状态不明的商品；
  若迷因热度退潮，存货风险完全由售卖端承担。
  读者在观察这条轨道时，应同时看到它的活力
  与它的脆弱——两者是同一个权利真空的两面。
\end{nlcritique}

\section{2026-03 事件的双重读法：法律层面与意义层面}\label{sec:09:split}

第四章留下了一个待展开的论证：
“断裂叙事”（奶蛙“正式脱离了奶龙”）与
“连续叙事”（授权、维权、关评的产业线）
不是对错关系，而是分别描述了\keyw{意义层面}
与\keyw{法律层面}的正典—变种关系。
现在给出这个论证。

\textbf{法律层面}：奶蛙与奶龙从未分离。
变异形象与正典形象的相似性正是侵权的构成基础
（也正因为相似，官方的拉黑才有明确对象）；
只要相似性存在，法律层面的支配关系
（权利人可以主张什么）就持续存在。
关评、拉黑、呼吁停止丑化，都是这一层面的动作
\conf{C}\cite{shoutanji2026naiwa}。

\textbf{意义层面}：奶蛙确实脱离了奶龙。
第六章的谱系分析显示，奶蛙的辨识单元
（黄色 + 大笑 + 膨腹）已经可以脱离正典设定独立运行；
“潘多拉奶盒”叙事是社群对这一脱离的自觉表述；
官方评论区“仍有无数奶蛙图片”
\conf{C}\cite{shoutanji2026naiwa}
则说明意义层面的流通不再服从法律层面的边界。

\begin{proposition}[双重读法（教学命题）]\label{prop:09:dual}
  2026-03 的官方抑制之所以“不彻底”，
  是因为抑制工具（关评、拉黑）只能作用于
  \keyw{法律层面可执行的对象}（官方账号评论区、
  特定账号），而奶蛙的存在方式已经转移到
  \keyw{意义层面不可执行的空间}（个体创作、
  群聊流通、线下立牌）。两个层面的错位
  构成“迷因生态的缝隙”：\keyw{失败的控制}
  不是控制不努力，而是控制对象与控制工具
  在层面上的错配。
\end{proposition}

命题 \ref{prop:09:dual}（\conf{M} 级，
事件基础为 \conf{C} 级）同时回答第三章的三条承诺：
其一，“官方控制 vs 迷因生态缝隙”的框架即本命题；
其二，“如果关评拉黑是彻底的，奶蛙现象在时间上会被截断”
——反事实推理的依据正是层面错配：
不存在能作用于意义层面的执行工具；
其三，受众外溢（第三章）的第三层
“民间再生产”之所以能独立完成“大学生化”，
正因为它在意义层面运行，不受内容内核层
与平台结构层的官方意图约束。

\section{收编：商业与迷因的双向运动}\label{sec:09:recoup}

Hebdige 描述了亚文化的经典命运：风格作为抵抗被
主流与商业\keyw{收编}（\en{incorporation}），
转化为商品与时尚\cite{hebdige1979subculture}。
奶蛙案例为这个理论提供了一个当代续集——
收编不是单行道。

\subsection{商业收编迷因}

民间周边（\ref{sec:09:variant}）是收编的第一形态：
把立场的符号（“反精致”的姿态，第七章 M3）
变成货架上的商品。亚文化理论预测，
收编会稀释风格的抵抗含义——
当奶蛙玩偶与正版 IP 玩偶并列销售时，
“转发一只丑蛙”的表态成本趋近于购买一只玩偶的成本。
第七章 M3 的“阶级自嘲”批判
（\ref{sec:07:m3} 的批判框）在这里得到经济侧印证：
\keyw{收编把表态变成了消费，把不合作变成了生意}。

\subsection{迷因反向收编商业}

奶蛙案例的特殊之处在于第二个方向：
\keyw{迷因逻辑也在改造商业行为}。
官方奶龙的六周年营销（五城大屏、单曲出道、
“歌手”人设\conf{B}\cite{chinadaily2026six}）
在形式上就是对迷因式流通的学习——
把 IP 做成“可传播性单位”投放；
2026-01-25 登上总台《动漫盛典》的“高姿态”
\conf{B}\cite{chinadaily2026six}
与同月内社群里奶蛙的刷屏，构成同一文化市场上
两种流通逻辑的并置。一篇 2026 年 7 月的评述
甚至以“奶龙，正在成为一件艺术品”为题，
讨论正典奶龙被“艺术化”“档案化”的趋势
\conf{C}\cite{yuhuoxingxing2026}——
无论该评述的判断是否成立，
“一个子供向 IP 被讨论为艺术品”这件事本身，
就是迷因文化抬升其讨论语域的证据。

\begin{nltermbox}{术语：收编（incorporation）}
  Hebdige 的亚文化理论术语：主流文化与商业把
  亚文化风格转化为可销售商品、可传播时尚，
  从而消解其抵抗性的过程\cite{hebdige1979subculture}。
  本书用法提醒：收编预设了“抵抗”的存在，
  而 M3 的判别问题（\ref{sec:07:m3} 批判框）
  尚未解决——若“反精致”本就是人设消费，
  则周边商品化不是“收编抵抗”，
  只是普通的供给跟随需求。
  \keyw{使用该术语前，请先确认抵抗性证据}。
\end{nltermbox}

\section{本章小结}\label{sec:09:summary}

本章在融合文化与可传播性的底图上，
对照了两条经济轨道：官方的“内容 + 授权”双轮
（维权是授权的前置条件；表情包 80 亿次说明官方
主动模仿迷因式流通但对失控流通保持警惕），
与民间的零授权可传播性经济（创作者不获酬的结构性
报酬真空）。对 2026-03 事件给出了法律—意义双重读法
（命题 \ref{prop:09:dual}）：抑制工具只能作用于
法律层面的可执行对象，意义层面的流通无法执行，
层面错配即“迷因生态的缝隙”。收编分析指出
商业收编迷因与迷因反向收编商业的双向运动。
遗留问题有二：其一，两条轨道在校园空间里如何交汇
（立牌同时是零授权符号商品与社群仪式物），
第十章处理；其二，本章对“意义层面”的分析
全部依赖社群文本与 \conf{C} 级材料，
意义脱离的实际程度需要受体研究提供行为证据
（第十一章）。

\begin{nlexercises}
  \item 用命题 \ref{prop:09:unit} 的“不依赖原始语境”
        三个维度（内容、署名、载体），
        分别检验“结尾集体大笑”“奶蛋”“某品牌联名奶茶杯”
        三个构件，排序其可传播性并说明依据。
  \item 案例 9-1 认为官方“不反对迷因式流通本身，
        反对失控的流通”。请为“失控”给出一个
        \keyw{可操作的判据}（什么条件下流通算失控），
        并讨论该判据分别从官方视角与社群视角
        写出来会有什么不同。
  \item 表 \ref{tab:09:two} 指出民间轨道中
        “创作者基本不获酬”。有观点认为
        “匿名创作本就放弃了报酬，无所谓真空”；
        另一观点认为“社群中人人都在免费使用，
        报酬通道的缺失是结构不公”。
        写出两种观点各自的隐含前提，
        并说明本章的 E3 规则（权利保留）支持哪一方的
        \keyw{分析姿态}、不支持哪一方的\keyw{结论}。
\end{nlexercises}

\nlkeywords 授权体系、内容 + 授权、可传播性、
可传播性单位、双轨分发、零授权经济、收编、
法律层面、意义层面、迷因生态缝隙、创作者报酬
